\hypertarget{boundary-and-assumptions}{%
\section{Boundary and assumptions}\label{boundary-and-assumptions}}

The Paxos participant is \texttt{KVReplica}, instantiated as one
homogeneous cluster of exactly three members: \texttt{kv-0},
\texttt{kv-1}, and \texttt{kv-2}. \texttt{Browser} and \texttt{WebApp}
are actors but are not Paxos participants. Two independent browser
sessions may have tasks in flight concurrently.

The confirmed availability premise means that at every modeled instant
at least two replicas are running, all replicas can reach one another
while running, and each can reach its durable storage. This gives a live
majority. Persistence additionally requires an acceptor to fsync its
accepted ballot, slot, key, and value before sending \texttt{Accepted};
the web response is not successful until a quorum of two durable
acceptances exists. A restarted replica recovers its accepted log before
participating. If all three replicas stop, the service is unavailable
rather than forgetful: after restart, quorum recovery reconstructs the
committed log from durable acceptor state before reads or writes resume.

The confirmed conflict rule is last-committed-slot-wins: if the two
browser sessions write the same key concurrently, the leader assigns
distinct Paxos log slots and the value in the later committed slot is
displayed.

\hypertarget{related-grammar-set}{%
\section{Related grammar set}\label{related-grammar-set}}

\begin{verbatim}
project paxos_kv ::=
  repeat(quorum_write || linearizable_list || recover_replica)

quorum_write ::=
  Browser.Put(key, value, session, request_id) ;
  WebApp.Propose(key, value, session, request_id) ;
  Leader.Assign(ballot, slot) ;
  parallel(
    Replica.Accept(ballot, slot, key, value) ;
    Replica.Accept(ballot, slot, key, value) ;
    Replica.Accept(ballot, slot, key, value)
  ) ;
  mOfN(2, Replica.DurableAccepted(ballot, slot)) ;
  Leader.Commit(ballot, slot, key, value) ;
  WebApp.PutSucceeded(key, value, slot) ;
  ref(linearizable_list)

linearizable_list ::=
  Browser.List(session, request_id) ;
  WebApp.ReadBarrier(ballot) ;
  mOfN(2, Replica.ReadBarrierAck(ballot, committed_slot)) ;
  Leader.Snapshot(committed_slot, entries) ;
  WebApp.ListSucceeded(committed_slot, entries)

recover_replica ::=
  Replica.Unavailable(instance) ;
  Replica.Restarted(instance) ;
  Replica.RecoverAcceptedLog(instance) ;
  Replica.CatchUp(committed_slot) ;
  Replica.Available(instance)
\end{verbatim}

The notation above is intentionally behavioral. An implementation may
use a stable Multi-Paxos leader, but leader election and ballot changes
must remain observable whenever leadership changes.

\hypertarget{two-session-interaction-projection}{%
\section{Two-session interaction
projection}\label{two-session-interaction-projection}}

This is the interaction structure to be generated from the grammar. The
two browser branches are incomparable until the leader assigns slots;
either may arrive first.

\begin{verbatim}
Browser A      Browser B        WebApp          KV leader       KV followers
    | PUT a=1      |               |                 |                 |
    |------------->|               |                 |                 |
    |               | PUT b=2      |                 |                 |
    |               |------------->|                 |                 |
    |               |               | propose A      |                 |
    |               |               |--------------->|                 |
    |               |               | propose B      |                 |
    |               |               |--------------->|                 |
    |               |               |          assign slots s,s+1      |
    |               |               |                 | accept/fsync -->|
    |               |               |                 |<-- accepted (2) |
    |               |               |<-- commit A ----|                 |
    |<-- success + current list ----|                 |                 |
    |               |               |                 | accept/fsync -->|
    |               |               |                 |<-- accepted (2) |
    |               |               |<-- commit B ----|                 |
    |               |<-- success + current list ------|                 |
\end{verbatim}

The durable event representation must keep separate
\texttt{(session,\ scenario)} keys and use \texttt{priorIds} to record
immediate predecessor order. \texttt{priorIds\ =\ {[}{]}} means that an
event is a minimal (first) event in the scenario. If several events have
no predecessor, they are mutually unordered and may proceed in parallel.
A singleton list is the ordinary prior-event link; a list with several
IDs is a join. These links are ordering backpointers, not claims of
causation. Event file line order must not serialize otherwise unordered
browser requests.

\hypertarget{observable-safety-and-progress-obligations}{%
\section{Observable safety and progress
obligations}\label{observable-safety-and-progress-obligations}}

\begin{itemize}
\tightlist
\item
  A successful PUT necessarily has two prior \texttt{DurableAccepted}
  observations for the same ballot, slot, key, and value.
\item
  No two different values are committed for the same log slot.
\item
  Every committed slot remains present after any one replica is
  unavailable.
\item
  A list result reflects one committed prefix selected after a quorum
  read barrier; it does not combine entries from different prefixes.
\item
  With two mutually reachable durable replicas and a leader eventually
  chosen, every accepted browser request eventually succeeds or returns
  an explicit error.
\item
  Replica unavailability, leadership loss, retry, rejection, and
  recovery are visible outcomes rather than hidden timeouts.
\item
  At most three replica instances exist, all implementing the same
  \texttt{KVReplica} actor specification.
\end{itemize}

\hypertarget{prompt-derived-outputs}{%
\section{Prompt-derived outputs}\label{prompt-derived-outputs}}

\begin{description}
\item[\texttt{paxos.write.interaction}] (interaction diagram) How does a browser write become persistent? \emph{Reducer:} Partition write events by \texttt{(session,request\_id)} and render quorum branches in \texttt{priorIds} order, with typed message fields in each node
\item[\texttt{paxos.write.state\_machine}] (state machine) Which observable write states and failure exits exist? \emph{Reducer:} Render residuals of \texttt{quorum\_write}
\item[\texttt{paxos.list.interaction}] (interaction diagram) How is the complete key/value listing made linearizable? \emph{Reducer:} Render list and read-barrier events by \texttt{(session,request\_id)}
\item[\texttt{paxos.list.state\_machine}] (state machine) Which observable list states and failure exits exist? \emph{Reducer:} Render residuals of \texttt{linearizable\_list}
\item[\texttt{paxos.concurrent.order}] (prior-pointer graph) Can two browser sessions access the app concurrently? \emph{Reducer:} Render incomparable request branches, slot assignment, quorum joins, and responses
\item[\texttt{paxos.commit\_latency}] (2D function) How does write latency change with concurrent browser sessions? \emph{Reducer:} Pair PUT start to success/error by prior ID; x=max concurrent sessions, y=latency milliseconds
\item[\texttt{paxos.quorum\_size}] (scalar and histogram) Did every acknowledged write have a durable majority? \emph{Reducer:} Count distinct accepting replica instances per committed slot
\item[\texttt{paxos.replica\_availability}] (line) Were at least two replicas available throughout the observation window? \emph{Reducer:} Sweep \texttt{Unavailable}/\texttt{Available} intervals and plot available replica count over monotonic time
\item[\texttt{paxos.recovery\_lag}] (2D function) How far behind is a restarting replica? \emph{Reducer:} x=time since restart; y=leader committed slot minus replica applied slot
\item[\texttt{paxos.outcomes}] (pie/histogram) What happened to browser PUT and list requests? \emph{Reducer:} Count success, explicit rejection, timeout/censored, and leadership-loss outcomes
\end{description}


Numeric reductions use these stable output IDs at \texttt{/metrics}.
Every diagram, state machine, and chart is derived from the accepted
grammar and events rather than maintained as a second specification.

\hypertarget{web-behavior}{%
\section{Web behavior}\label{web-behavior}}

The page contains a string \texttt{key} input, a string \texttt{value}
input, and a submit action. Below the form it renders the linearizable
listing of all committed key/value pairs. After a successful write the
initiating session receives a page containing a listing at least as new
as its committed write. The other browser session observes that write on
its next list operation; server push is not currently required.

There is no authentication and no delete operation. Keys are nonempty
UTF-8 strings of at most 128 encoded bytes. Values are UTF-8 strings of
at most 4096 encoded bytes. The database holds at most 10,000 distinct
keys and supports at least two browser sessions with requests in flight
concurrently. Invalid or capacity-exceeding writes return explicit
errors without changing the committed log. Each successful submission
returns a listing containing at least that committed write.
